\begingroup
\renewcommand{\arraystretch}{1.45}
\fancyhead[L]{}
\fancyhead[C]{\small Das wunderbare Pentagramm}
\par\medskip\Needspace{5\baselineskip}\pdfbookmark[1]{Art. 6}{pentagramma-art-6}\centerline{[6].}\nobreak\smallskip
\noindent\hspace*{2em}$A,\,B,\,C,\,D,\,E$\quad die Eckpunkte des Polygons\par
\noindent\hspace*{2em}$a,\quad b,\quad c,\quad d,\quad e$\quad Pole der Diagonalen\par
\noindent\hspace*{2em}$0,\quad1,\quad2$\quad die drei Hauptachsen entsprechend den Wurzeln
$G,G',G''$ der Gleichung
\[
\begin{gathered}\frac{t(2t-1)^2}{t-1}\\[9pt]\quad{}=(\operatorname{tang}AB\mathord{.}\operatorname{tang}BC\mathord{.}\operatorname{tang}CD\mathord{.}\operatorname{tang}DE\mathord{.}\operatorname{tang}EA)^2\end{gathered}
\]
wo für $G$ die negative Wurzel genommen werden mag; sodass
\[
Guu++G'u'u'+G''u''u''=0
\]
wenn $u$ die Koordinaten irgendeines der Punkte $A,B,C,D,E$ bedeuten,
also zugleich $uu+u'u'+u''u''=1$.


\Needspace{18\baselineskip}
Die Quelle der Hauptsätze ist in den zwei Gleichungen enthalten
\[
\begin{gathered} \text{I.}\qquad\cos0A\mathord{.}\cos0b=-\frac{\operatorname{tang}EA\mathord{.}(2G-1-\operatorname{tang}AB^2)}{4(G-G')(G-G'')}\\[9pt] \text{II.}\qquad\cos0A\mathord{.}\cos0C=-\frac{2G-1-\frac{1}{\sin DE^2}}{4\cos BC\mathord{.}\cos AB(G-G')(G-G'')} \end{gathered}
\]


\Needspace{44\baselineskip}
\SourceLeaf{498}

Die Gleichung I. repräsentiert 30 Gleichungen, die II. hingegen 15, da alle
Permutationen der Achsen und der Eckpunkte erlaubt sind. Noch zierlicher
(in den anfänglichen Bezeichnungen, wo $\alpha=\operatorname{tang}CD^2$ usw.)
\begingroup\small\setlength{\arraycolsep}{3pt}\[
\resizebox{.95\linewidth}{!}{$\displaystyle \begin{array}{ll} \mathrm{I.}&\left\{\begin{aligned} \cos 0C\mathbin{.}\cos 0d&=\frac{1+\alpha-2G}{4(G'-G)(G''-G)}\mathbin{.}\surd\varepsilon=\mathfrak{A}\mathbin{.}\surd\varepsilon\\[9pt] \cos 0D\mathbin{.}\cos 0c&=\frac{1+\alpha-2G}{4(G'-G)(G''-G)}\mathbin{.}\surd\beta=\mathfrak{A}\mathbin{.}\surd\beta \end{aligned}\right.\\[18pt] \mathrm{II.}&\displaystyle \cos 0B\mathbin{.}\cos 0E=\frac{2\alpha+1-2\alpha G}{4(G'-G)(G''-G)}\mathbin{.}\surd\beta\varepsilon=\mathfrak{a}\mathbin{.}\surd\beta\varepsilon \end{array}$}
\]\endgroup

\Needspace{28\baselineskip}
Die zehn Gleichungen I. können nur für neun gelten, weil die Multiplikation von fünfen dasselbe Resultat gibt wie die Multiplikation der fünf übrigen.
Es muss also zwischen den 10 Größen $\mathfrak{A},\mathfrak{a},\mathfrak{B},\mathfrak{b}$ usw. vier Bedingungsgleichungen geben, welche am zierlichsten so dargestellt werden
\[
\begin{gathered} \beta\mathfrak{b}\mathfrak{c}\mathfrak{C}=\varepsilon\mathfrak{e}\mathfrak{d}\mathfrak{D},\\[9pt] \gamma\mathfrak{c}\mathfrak{d}\mathfrak{D}=\alpha\mathfrak{a}\mathfrak{e}\mathfrak{E}\quad\text{usw. oder auch}\\[9pt] \beta\mathfrak{a}\mathfrak{A}\mathfrak{C}=\gamma\mathfrak{d}\mathfrak{B}\mathfrak{D},\\[9pt] \gamma\mathfrak{b}\mathfrak{B}\mathfrak{D}=\delta\mathfrak{e}\mathfrak{C}\mathfrak{E}\quad\text{usw.} \end{gathered}
\]
Es ist aber
\[
\begin{gathered}\cos 0A^2=\frac{\mathfrak{e}\mathfrak{a}\mathfrak{b}}{\mathfrak{c}\mathfrak{d}}\mathbin{.}\alpha =\frac{\mathfrak{D}\mathfrak{b}\gamma}{\mathfrak{E}} =\frac{\mathfrak{C}\mathfrak{c}\delta}{\mathfrak{B}}\quad\text{usw.,} \\[9pt]\quad{}\qquad\cos 0a^2=\frac{\mathfrak{C}\mathfrak{D}}{\mathfrak{a}}\quad\text{usw.}\end{gathered}
\]

\clearpage
\endgroup
